\documentclass[UTF8,12pt]{ctexart}
\usepackage[a4paper,margin=24mm]{geometry}
\usepackage{amsmath,amssymb,bm,graphicx,adjustbox}
\newcommand{\fitmath}[1]{\begin{adjustbox}{max width=\linewidth}$\displaystyle #1$\end{adjustbox}}
\setlength{\parindent}{0pt}
\setlength{\parskip}{.7em}
\sloppy
\allowdisplaybreaks
\begin{document}
\section*{2012 年考研数学三 · 参考答案}
\subsection*{原 PDF 第 40 页}

\subsection*{2012 年真题参考答案}

\subsection*{一、选择题}

(1) C。 (2) A。 (3) B。 (4) D。 (5) C。 (6) B。 (7) D。 (8) B。


\subsection*{二、填空题}

(9) \(\displaystyle e^{-\sqrt2}\)。 (10) \(\displaystyle \dfrac{\displaystyle 1}{\displaystyle e}\)。 (11) \(\displaystyle 2dx-dy\)。 (12) \(\displaystyle 4\ln2\)。 (13) \(\displaystyle -27\)。 (14) \(\displaystyle \dfrac{\displaystyle 3}{\displaystyle 4}\)。


\subsection*{三、解答题}

(15) \(\displaystyle \dfrac{\displaystyle 1}{\displaystyle 12}\)。


(16) \(\displaystyle \dfrac{\displaystyle 1}{\displaystyle 2}\)。


(17)（I）\(\displaystyle C(x,\allowbreak y)=20x+\dfrac{\displaystyle x^2}{\displaystyle 4}+6y+\dfrac{\displaystyle y^2}{\displaystyle 2}+10000\)。（II）若总产量为 50 件，则当甲种产品的产量为 24 件，乙种产品的产量为 26 件时，总成本最小，最小成本为 \(\displaystyle C(24,\allowbreak 26)=11118\) 万元。（III）当总产量为 50 件且总成本最小时，甲种产品的边际成本为 32 万元。其经济意义为：当甲、乙两种产品的产量分别为 24、26 件时，若甲种产品的产量再增加 1 件，则总成本增加 32 万元，即当乙种产品的产量为 26 件时，生产第 25 件甲种产品需要 32 万元。


(18) 证明略。


(19)（I）\(\displaystyle f(x)=e^x\)。（II）\(\displaystyle (0,\allowbreak 0)\)。


(20)（I）\(\displaystyle |A|=1-a^4\)。（II）\(\displaystyle a=-1\) 时，方程组 \(\displaystyle Ax=\beta\) 有无穷多解，其通解为 \(\displaystyle k(1,\allowbreak 1,\allowbreak 1,\allowbreak 1)^{\mathrm T}+(0,\allowbreak -1,\allowbreak 0,\allowbreak 0)^{\mathrm T}\)，其中 \(\displaystyle k\) 为任意常数。


(21)（I）\(\displaystyle a=-1\)。（II）\(\displaystyle Q=\begin{pmatrix}\dfrac{\displaystyle 1}{\displaystyle \sqrt6}&\dfrac{\displaystyle 1}{\displaystyle \sqrt2}&\dfrac{\displaystyle 1}{\displaystyle \sqrt3}\\\dfrac{\displaystyle 1}{\displaystyle \sqrt6}&-\dfrac{\displaystyle 1}{\displaystyle \sqrt2}&\dfrac{\displaystyle 1}{\displaystyle \sqrt3}\\\dfrac{\displaystyle 2}{\displaystyle \sqrt6}&0&-\dfrac{\displaystyle 1}{\displaystyle \sqrt3}\end{pmatrix}\)，正交变换 \(\displaystyle x=Qy\) 将二次型 \(\displaystyle f(x_1,\allowbreak x_2,\allowbreak x_3)\) 变成标准形 \(\displaystyle f=6y_1^2+2y_2^2\)。


(22)（I）\(\displaystyle \dfrac{\displaystyle 1}{\displaystyle 4}\)。（II）\(\displaystyle -\dfrac{\displaystyle 2}{\displaystyle 3}\)。


(23)（I）\(\displaystyle f_V(v)=\begin{cases}2e^{-2v},\allowbreak &v>0,\allowbreak \\0,\allowbreak &v\leq0.\end{cases}\)（II）\(\displaystyle 2\)。


\clearpage
\end{document}
